\chapter{Architecture backend (FastAPI)}
\label{chap:backend}

\section{Organisation}
Le backend est structuré en \textbf{routeurs spécialisés} (un par domaine fonctionnel),
montés sous le préfixe \texttt{/api}. Au démarrage, l'application initialise la base de
données, charge l'ensemble des modèles via le \texttt{ModelRegistry}, instancie la flotte
de démonstration et active les connecteurs industriels éventuels. Le
tableau~\ref{tab:routes} recense les principaux groupes de routes.

\begin{table}[h]
\centering\small
\caption{Principaux groupes de routes de l'API.}
\label{tab:routes}
\begin{tabular}{>{\bfseries}l p{9.5cm}}
\toprule
Domaine & Routes principales \\
\midrule
Prédiction & \texttt{/predict}, \texttt{/model/select}, \texttt{/model/active/\{ds\}},
             \texttt{/model/versions/\{ds\}/\{m\}}, \texttt{/model/rollback} \\
Flotte / santé & \texttt{/fleet}, \texttt{/machine/\{id\}}, \texttt{/machine/\{id\}/history} \\
Simulation & \texttt{/ws/simulation} (WebSocket), \texttt{/simulation/units} \\
Analyse spectrale & \texttt{/signal/indicators}, \texttt{/signal/spectrum},
             \texttt{/signal/diagnose}, \texttt{/signal/iso-severity}, \texttt{/bearing/frequencies} \\
Ingestion & \texttt{/ingest/signal}, \texttt{/ingest/baseline/\{finalize,reset,status\}}, \texttt{/replay/\{ds\}/\{idx\}} \\
Anomalie & \texttt{/anomaly/ensemble/\{ds\}}, \texttt{/anomaly/ensemble/demo/\{ds\}} \\
KPIs / analytics & \texttt{/kpi/real/\{id\}}, \texttt{/kpi/roi/\{id\}},
             \texttt{/analytics/prediction-accuracy/\{id\}}, \texttt{/analytics/false-alarm-rate} \\
GMAO & \texttt{/workorder(s)}, \texttt{/maintenance/event} \\
Confiance & \texttt{/drift/status}, \texttt{/audit/recent}, \texttt{/explain/shap},
             \texttt{/explain/calibration/\{ds\}} \\
Copilot & \texttt{/copilot/chat}, \texttt{/copilot/report/\{id\}} \\
Sécurité & \texttt{/auth/register}, \texttt{/auth/login}, \texttt{/auth/me} \\
Retrain & \texttt{/retrain/start}, \texttt{/retrain/demo}, \texttt{/retrain/status} \\
\bottomrule
\end{tabular}
\end{table}

\section{Le registre de modèles}
Le \texttt{ModelRegistry} est la pièce maîtresse du backend. Il centralise le chargement
(registre benchmark, auto-encodeurs, CNN, MLP spécialisés), la gestion du modèle actif par
dataset (persistée en JSON), l'inférence unifiée, le réentraînement asynchrone et le
versioning. L'interface de prédiction unifiée applique systématiquement la même logique :
sélection du modèle actif (ou imposé), normalisation conditionnelle des entrées (uniquement
pour le MLP et le Huber, sensibles à l'échelle), exécution du monitoring de dérive en
parallèle, puis production d'un résultat au format homogène. Pour une tâche de régression
(CMAPSS), la sortie est la RUL moyenne plafonnée à 125 cycles ; pour une tâche de
classification, elle comprend le label décodé (via le \emph{LabelEncoder} du dataset) et un
niveau de confiance dérivé des probabilités. Ce format uniforme, indépendant du modèle
sous-jacent, simplifie la consommation côté client et l'intégration de nouveaux modèles.

\section{La simulation temps réel}
Le c\oe ur dynamique de la plateforme est le WebSocket de simulation. À chaque cycle, pour
la machine surveillée, le serveur : (i) sélectionne la fenêtre de données réelle suivante,
(ii) exécute la prédiction supervisée et le score d'anomalie, (iii) calcule l'indice de
santé lissé, (iv) reconstruit un signal exploitable pour l'affichage spectral, (v) met à
jour le \texttt{HealthTracker} (historique, tendance, alertes) et (vi) diffuse un message
JSON complet. Toutes les cinq itérations, l'événement est consigné dans le journal d'audit
et persisté en base. Un \texttt{ConnectionManager} isole l'état de chaque client (machine,
dataset, vitesse), permettant la surveillance simultanée de plusieurs machines.

\section{Sécurité et persistance}
L'\textbf{authentification} repose sur des jetons \textbf{JWT} signés (HS256, expiration
configurable) et des mots de passe hachés en \textbf{bcrypt}. Deux rôles sont gérés :
\texttt{admin} (réentraînement, injection de défauts, rollback) et \texttt{user}
(consultation). La \textbf{persistance} s'appuie sur SQLAlchemy : les modèles
\texttt{MachineState}, \texttt{Alert}, \texttt{DiagnosticLog}, \texttt{ModelMetadata},
\texttt{User}, \texttt{MaintenanceEvent}, \texttt{ModelVersion}, \texttt{WorkOrder} et
\texttt{AlertVerdict} couvrent l'ensemble du domaine. Le mode \textbf{WAL} de SQLite autorise
lectures et écriture concurrentes, indispensable à la coexistence de la boucle de simulation,
des écritures d'événements et du versioning.

\section{Validation des entrées}
Toutes les entrées sont validées par des schémas \textbf{Pydantic} (longueur minimale de
signal, positivité des fréquences, géométrie de roulement cohérente), ce qui sécurise l'API
contre les données malformées et documente automatiquement le contrat d'interface via
Swagger (\texttt{/docs}).
